\section{The Huntington-Hill Precedent}

Congressional redistricting is not the first time American democracy has confronted the need for algorithmic objectivity. For 150 years, Congress struggled with an analogous problem: apportioning House seats among states based on decennial census counts. From 1790 to 1941, every census triggered political battles over which mathematical method to use, with each formula favoring different states. Congress changed the apportionment method six times, always to advantage particular political coalitions~\cite{balinski1982fair}.

In 1941, Congress adopted the Huntington-Hill method of equal proportions, a
formula that has supplied a durable operational allocation rule for eight
decades~\cite{huntington1928methods}. The method did not end disputes over
House size, census counts, constitutional interpretation, or alternative
apportionment axioms. Its narrower achievement was to remove cycle-specific
allocation discretion after Congress selected the method and inputs.

This precedent suggests a pathway for redistricting. Just as apportionment answered ``How many seats per state?'' with mathematical objectivity, redistricting might answer ``Where to draw district boundaries?'' through algorithmic methods. The parallel is instructive:

\begin{center}
\begin{tabular}{ll}
\textbf{Apportionment (operational rule)} & \textbf{Redistricting (open rule)} \\
How many seats per state? & Where to draw boundaries? \\
Mathematical formula (1941) & Partisan negotiation \\
Durable post-enactment procedure & Persistent method and boundary conflict \\
\end{tabular}
\end{center}

The relevant question is whether Congress can create a similarly durable
post-enactment decision procedure for a harder problem while preserving legal
and representational judgment.
